\begin{pastexam}[label={px:nchu-115-is-15}]{中興資工所乙組 115 學年度 第 15 題}
（單選；每題 2 分，答錯不倒扣）

You are designing an auto-complete feature for a mobile keyboard app. The system must
efficiently return all dictionary words that share a given prefix, such as ``com--''. What data
structure is most suitable for storing the dictionary?
\begin{enumerate}[label=\Alph*.,nosep]
\item A binary search tree storing words in lexicographical order, enabling in-order traversal
for sorted output.
\item A hash set storing all full words and scanning for any entry that begins with the requested
prefix.
\item A Trie (prefix tree) where each character corresponds to a node, allowing fast
prefix-based lookup.
\item A balanced heap storing all words based on their length, allowing quick access to short
words.
\end{enumerate}
\end{pastexam}

\begin{pastsoln}
答案是 C。

Trie 從根往下，每一層對應字的一個字元（資結書第~11.5.1~節）。 有相同前綴的字，在 trie 裡共用
同一條從根出發的路徑。 查前綴 \texttt{com}：從根照 \texttt{c}、\texttt{o}、\texttt{m} 走三步，
走到的那個節點底下的子樹就是所有以 \texttt{com} 開頭的字，把子樹走一遍全部列出。 找到節點只要
$O(\text{前綴長度})$，和字典有多少字無關；列出答案的時間和答案的總長度成正比。
\begin{itemize}
\item A：以字典序排的 BST 其實也做得到，先找第一個 $\ge$ \texttt{com} 的字，再依中序往後走到
不再以 \texttt{com} 開頭為止。 但每一步比較都是整個字串的比較，平衡時還要 $O(\log n)$ 次，
而且題目描述的用途是「排序輸出」，不是前綴查詢。 和 trie 比，不是最合適的（資結書
第~11.5.4~節）。
\item B 錯：雜湊表只支援「整個字一樣」的查詢，雜湊值和前綴沒有關係，只能把 $n$ 個字全部掃過
一遍，$O(n)$。
\item C 對。
\item D 錯：堆積只能快速拿到最大（或最小）的一個，依長度排和前綴毫無關係。
\end{itemize}
\end{pastsoln}
